% GENERATED FILE. Edit canonical lessons, then run npm run generate:books.
% canonical-source-hash: fnv1a64:3b7d263aca1a3449
% canonical-lessons: JA-C376-shinkansen, JA-C376-kaisoku, JA-C376-shitetsu, JA-C376-rosen, JA-C376-tobasu

\chapter{Fast Trains and Local Lines}
\label{ch:ja-fast-trains-and-local-lines}
\hlchaptermodality{376}

\begin{chapteropening}
\textbf{By the end of this chapter:} \emph{I can choose between fast trains and local lines.}

Complete the last lesson of chapter 376: I can choose between fast trains and local lines.
\end{chapteropening}

\section[\texorpdfstring{\ja{しんかんせん} (shinkansen)}{しんかんせん (shinkansen)}]{\ja{しんかんせん} (shinkansen) — the bullet train}

\label{lesson:JA-C376-shinkansen}



\begin{quote}
Before the new one: say the Japanese for to ride past one's stop, then the Japanese for to get on board.
\end{quote}

\subsection*{You'll want to know: \ja{しんかんせん}}
\textbf{\ja{しんかんせん}} — \emph{shinkansen} — \textquotedblleft{}the bullet train\textquotedblright{}.

\begin{quote}
\textbf{\ja{しんかんせんは} \ja{ひこうきより} \ja{べんりです。}} — \emph{shinkansen wa hikōki yori benri desu} — \textquotedblleft{}The bullet train is more convenient than the plane.\textquotedblright{}
\end{quote}

\begin{cousinweb}
\textbf{\ja{しんかんせん}} (\emph{shinkansen}) is Sino-Japanese (kango): three Chinese characters read with their on'yomi, readings borrowed from Chinese. They are \emph{shin}, \textquotedblleft{}new\textquotedblright{} (Middle Chinese \emph{sin}), \emph{kan}, \textquotedblleft{}main part\textquotedblright{} (\emph{kanH}) and \emph{sen}, \textquotedblleft{}line\textquotedblright{} (\emph{sjenH}) (KANJIDIC2; Middle Chinese in Baxter's transcription).
\end{cousinweb}

\begin{grammarlens}[title={What you've built}]
One more word for getting around.
\end{grammarlens}

\subsection*{Your turn}
\begin{itemize}
\raggedright
  \item \emph{Say it:} \emph{shinkansen}
  \item \emph{Say it:} \emph{shinkansen}, once more
  \item \emph{Say it:} \emph{shinkansen wa hikōki yori benri desu} — The bullet train is more convenient than the plane.
  \item \emph{Recall:} say the Japanese for to ride past one's stop, then the Japanese for to get on board, then say \emph{shinkansen} again
\end{itemize}

\begin{tcolorbox}[breakable,colback=teal!4,colframe=teal!35!black,title={Before you move on}]
What does \ja{しんかんせん} mean? (\textquotedblleft{}The bullet train\textquotedblright{}.) Say it once more.
\end{tcolorbox}
\section[\texorpdfstring{\ja{かいそく} (kaisoku)}{かいそく (kaisoku)}]{\ja{かいそく} (kaisoku) — a rapid train}

\label{lesson:JA-C376-kaisoku}



\begin{quote}
Before the new one: say the Japanese for to get on board, then the Japanese for the bullet train.
\end{quote}

\subsection*{You'll want to know: \ja{かいそく}}
\textbf{\ja{かいそく}} — \emph{kaisoku} — \textquotedblleft{}a rapid train\textquotedblright{}.

\begin{quote}
\textbf{\ja{かいそくは} \ja{ちいさい} \ja{えきに} \ja{とまりません。}} — \emph{kaisoku wa chiisai eki ni tomarimasen} — \textquotedblleft{}The rapid train does not stop at small stations.\textquotedblright{}
\end{quote}

\begin{cousinweb}
\textbf{\ja{かいそく}} (\emph{kaisoku}) is Sino-Japanese (kango): two Chinese characters read with their on'yomi, readings borrowed from Chinese. They are \emph{kai}, \textquotedblleft{}pleasant\textquotedblright{} (Middle Chinese \emph{khwaejH}) and \emph{soku}, \textquotedblleft{}fast\textquotedblright{} (\emph{suwk}) (KANJIDIC2; Middle Chinese in Baxter's transcription).
\end{cousinweb}

\begin{grammarlens}[title={What you've built}]
One more word for getting around.
\end{grammarlens}

\subsection*{Your turn}
\begin{itemize}
\raggedright
  \item \emph{Say it:} \emph{kaisoku}
  \item \emph{Say it:} \emph{kaisoku}, once more
  \item \emph{Say it:} \emph{kaisoku wa chiisai eki ni tomarimasen} — The rapid train does not stop at small stations.
  \item \emph{Recall:} say the Japanese for to get on board, then the Japanese for the bullet train, then say \emph{kaisoku} again
\end{itemize}

\begin{tcolorbox}[breakable,colback=teal!4,colframe=teal!35!black,title={Before you move on}]
What does \ja{かいそく} mean? (\textquotedblleft{}A rapid train\textquotedblright{}.) Say it once more.
\end{tcolorbox}
\section[\texorpdfstring{\ja{してつ} (shitetsu)}{してつ (shitetsu)}]{\ja{してつ} (shitetsu) — a private railway}

\label{lesson:JA-C376-shitetsu}



\begin{quote}
Before the new one: say the Japanese for the bullet train, then the Japanese for a rapid train.
\end{quote}

\subsection*{You'll want to know: \ja{してつ}}
\textbf{\ja{してつ}} — \emph{shitetsu} — \textquotedblleft{}a private railway\textquotedblright{}.

\begin{quote}
\textbf{\ja{この} \ja{まちには} \ja{してつの} \ja{えきが} \ja{ふたつ} \ja{あります。}} — \emph{kono machi ni wa shitetsu no eki ga futatsu arimasu} — \textquotedblleft{}This town has two stations on the private railway.\textquotedblright{}
\end{quote}

\begin{cousinweb}
\textbf{\ja{してつ}} (\emph{shitetsu}) is Sino-Japanese (kango): two Chinese characters read with their on'yomi, readings borrowed from Chinese. They are \emph{shi}, \textquotedblleft{}private\textquotedblright{} (Middle Chinese \emph{sij}) and \emph{tetsu}, \textquotedblleft{}iron\textquotedblright{} (\emph{thet}) (KANJIDIC2; Middle Chinese in Baxter's transcription).
\end{cousinweb}

\begin{grammarlens}[title={What you've built}]
One more word for getting around.
\end{grammarlens}

\subsection*{Your turn}
\begin{itemize}
\raggedright
  \item \emph{Say it:} \emph{shitetsu}
  \item \emph{Say it:} \emph{shitetsu}, once more
  \item \emph{Say it:} \emph{kono machi ni wa shitetsu no eki ga futatsu arimasu} — This town has two stations on the private railway.
  \item \emph{Recall:} say the Japanese for the bullet train, then the Japanese for a rapid train, then say \emph{shitetsu} again
\end{itemize}

\begin{tcolorbox}[breakable,colback=teal!4,colframe=teal!35!black,title={Before you move on}]
What does \ja{してつ} mean? (\textquotedblleft{}A private railway\textquotedblright{}.) Say it once more.
\end{tcolorbox}
\section[\texorpdfstring{\ja{ろせん} (rosen)}{ろせん (rosen)}]{\ja{ろせん} (rosen) — a route; a line (of trains or buses)}

\label{lesson:JA-C376-rosen}



\begin{quote}
Before the new one: say the Japanese for a rapid train, then the Japanese for a private railway.
\end{quote}

\subsection*{You'll want to know: \ja{ろせん}}
\textbf{\ja{ろせん}} — \emph{rosen} — \textquotedblleft{}a route; a line (of trains or buses)\textquotedblright{}.

\begin{quote}
\textbf{\ja{ちかてつの} \ja{ろせんは} \ja{ふくざつです。}} — \emph{chikatetsu no rosen wa fukuzatsu desu} — \textquotedblleft{}The subway routes are complicated.\textquotedblright{}
\end{quote}

\begin{cousinweb}
\textbf{\ja{ろせん}} (\emph{rosen}) is Sino-Japanese (kango): two Chinese characters read with their on'yomi, readings borrowed from Chinese. They are \emph{ro}, \textquotedblleft{}route\textquotedblright{} (Middle Chinese \emph{luH}) and \emph{sen}, \textquotedblleft{}line\textquotedblright{} (\emph{sjenH}) (KANJIDIC2; Middle Chinese in Baxter's transcription).
\end{cousinweb}

\begin{grammarlens}[title={What you've built}]
One more word for getting around.
\end{grammarlens}

\subsection*{Your turn}
\begin{itemize}
\raggedright
  \item \emph{Say it:} \emph{rosen}
  \item \emph{Say it:} \emph{rosen}, once more
  \item \emph{Say it:} \emph{chikatetsu no rosen wa fukuzatsu desu} — The subway routes are complicated.
  \item \emph{Recall:} say the Japanese for a rapid train, then the Japanese for a private railway, then say \emph{rosen} again
\end{itemize}

\begin{tcolorbox}[breakable,colback=teal!4,colframe=teal!35!black,title={Before you move on}]
What does \ja{ろせん} mean? (\textquotedblleft{}A route; a line (of trains or buses)\textquotedblright{}.) Say it once more.
\end{tcolorbox}
\section[\texorpdfstring{\ja{とばす} (tobasu)}{とばす (tobasu)}]{\ja{とばす} (tobasu) — to skip; to let fly}

\label{lesson:JA-C376-tobasu}



\begin{quote}
Before the new one: say the Japanese for a private railway, then the Japanese for a route.
\end{quote}

\subsection*{You'll want to know: \ja{とばす}}
\textbf{\ja{とばす}} — \emph{tobasu} — \textquotedblleft{}to skip; to let fly\textquotedblright{}.

\begin{quote}
\textbf{\ja{とっきゅうは} \ja{ちいさい} \ja{えきを} \ja{とばします。}} — \emph{tokkyū wa chiisai eki o tobashimasu} — \textquotedblleft{}The limited express skips the small stations.\textquotedblright{}
\end{quote}

Polite \textbf{\ja{とばします}}, \emph{tobashimasu}; a godan verb. Past (polite): \textbf{\ja{きゅうこうが} \ja{えきを} \ja{とばしました。}}, \emph{kyūkō ga eki o tobashimashita}.

\begin{cousinweb}
\textbf{\ja{とばす}} (\emph{tobasu}) is a native Japanese word (wago), inherited rather than borrowed from Chinese; its kanji was attached later as a kun'yomi (JMdict). It is the partner of \emph{tobu}, \textquotedblleft{}to fly\textquotedblright{}, that makes something fly or skip (JMdict).
\end{cousinweb}

\begin{grammarlens}[title={What you've built}]
One more word for getting around.
\end{grammarlens}

\subsection*{Your turn}
\begin{itemize}
\raggedright
  \item \emph{Say it:} \emph{tobasu}
  \item \emph{Say it:} \emph{tobasu}, once more
  \item \emph{Say it:} \emph{tokkyū wa chiisai eki o tobashimasu} — The limited express skips the small stations.
  \item \emph{Recall:} say the Japanese for a private railway, then the Japanese for a route, then say \emph{tobasu} again
\end{itemize}

\begin{tcolorbox}[breakable,colback=teal!4,colframe=teal!35!black,title={Before you move on}]
What does \ja{とばす} mean? (\textquotedblleft{}To skip; to let fly\textquotedblright{}.) Say it once more.
\end{tcolorbox}
